\documentclass[12pt]{article}
\usepackage[T1]{fontenc}
\usepackage[utf8]{inputenc}
\usepackage{lmodern}
\usepackage{textcomp}
\usepackage{enumitem}
\usepackage{geometry}
\geometry{margin=1in}
\begin{document}
\begin{center}
Answer Key - History - Graduate Level Assessment 15 \
\small (Answers are ordered exactly as in the question paper.)
\end{center}

\section*{Multiple Choice}
\begin{enumerate}[label=\arabic*.]
\item \textbf{Answer:} A
\\ \textit{Options:} A) Australopithecus afarensis; B) Homo naledi; C) Homo neanderthalensis; D) Neanderthals; E) Homo floresiensis
\\ \textit{Note:} The correct answer is Australopithecus afarensis because fossil evidence, including well-dated remains and notable specimens like Lucy, confirms that this hominid existed around 3.9 to 2.9 million years ago, with significant archaeological findings suggesting a timeline leading up to 1.8 million years ago when early human ancestors were evolving.
\item \textbf{Answer:} A
\\ \textit{Options:} A) 600,000 to 300,000 years ago; B) 100,000 years ago to present; C) 400,000 years ago to present; D) 800,000 to 10,000 years ago; E) 500,000 to 200,000 years ago
\\ \textit{Note:} The correct answer is supported by fossil and archaeological evidence indicating that anatomically modern Homo sapiens emerged in Africa between 600,000 and 300,000 years ago, marking a significant period in human evolution.
\item \textbf{Answer:} A
\\ \textit{Options:} A) in the Mesolithic, goods traveled less far and fewer goods were exchanged; B) in the Mesolithic, goods traveled farther but the same number of goods were exchanged; C) in the Mesolithic, goods traveled less far but the same number of goods were exchanged; D) in the Mesolithic, goods traveled the same distance but more goods were exchanged; E) in the Mesolithic, goods traveled farther but fewer goods were exchanged
\\ \textit{Note:} The correct answer highlights that during the Mesolithic period, trade networks were less expansive and focused on localized exchanges, contrasting with the Upper Paleolithic's broader, more extensive trade patterns that facilitated the exchange of a greater variety of goods over longer distances.
\item \textbf{Answer:} D
\\ \textit{Options:} A) Java Trench; B) Sunda Passage; C) Australia Trench; D) Wallace Trench; E) New Guinea Strait
\\ \textit{Note:} The Wallace Trench is the correct answer because it is a significant geological feature that marks the boundary between the Australian and Eurasian tectonic plates, located between New Guinea/Australia and Java/Borneo.
\item \textbf{Answer:} A
\\ \textit{Options:} A) advanced metallurgy.; B) monumental works.; C) a system of money.; D) maritime navigation technology.; E) complex religious systems.
\\ \textit{Note:} The correct answer is advanced metallurgy because ancient South American civilizations, such as the Inca, primarily utilized gold, silver, and copper without developing the high-temperature techniques necessary for ironworking, which limited their metallurgical advancements compared to contemporaneous civilizations in other regions.
\end{enumerate}

\section*{True / False}
\begin{enumerate}[label=\arabic*.]
\setcounter{enumi}{5}
\item \textbf{Answer:} True
\\ \textit{Options:} A) True; B) False
\\ \textit{Note:} According to the passage: Neo-plasticism
\item \textbf{Answer:} False
\\ \textit{Options:} A) True; B) False
\\ \textit{Note:} The correct answer is: Antarctica
\item \textbf{Answer:} False
\\ \textit{Options:} A) True; B) False
\\ \textit{Note:} The correct answer is: the critical discussion of myths
\item \textbf{Answer:} False
\\ \textit{Options:} A) True; B) False
\\ \textit{Note:} The correct answer is: Fars Province
\item \textbf{Answer:} False
\\ \textit{Options:} A) True; B) False
\\ \textit{Note:} The correct answer is: mockingbird
\end{enumerate}

\section*{Long Form Response}
\begin{enumerate}[label=\arabic*.]
\setcounter{enumi}{10}
\item \textbf{Answer:} Gamal Abdul Nasser
\\ \textit{Note:} Expected answer: Gamal Abdul Nasser
\item \textbf{Answer:} Surrounding the medieval core there is a ring of late 19 th- and early 20 th-century neighbourhoods, with newer neighbourhoods positioned farther out
\\ \textit{Note:} Expected answer: Surrounding the medieval core there is a ring of late 19 th- and early 20 th-century neighbourhoods, with newer neighbourhoods positioned farther out
\end{enumerate}

\vfill
\noindent\textit{End of Answer Key}
\end{document}